% Job posting: https://ca.linkedin.com/jobs/view/machine-learning-engineer-at-invision-ai-4460322410
% =====================================================================
%  CV - Nishal Stanislaus Silva, PhD
%  Target: Invision AI - Machine Learning Engineer, Toronto ON (hybrid)
%  Variant: V5 (Computer Vision / Industrial)
% =====================================================================

\documentclass[11pt]{article}
\usepackage[left=0.7in,top=0.7in,right=0.7in,bottom=0.7in]{geometry}
\usepackage{enumitem}
\usepackage[hidelinks]{hyperref}
\usepackage{titlesec}
\usepackage{array}
\usepackage{setspace}
\usepackage{needspace}
\usepackage[T1]{fontenc}
\usepackage{newtxtext,newtxmath}
\ifdefined\pdfgentounicode
  \input{glyphtounicode}
  \pdfgentounicode=1
\fi
\setstretch{1.05}
\pagenumbering{gobble}
\titleformat{\section}{\Large\scshape}{}{4em}{}[\vspace{-0.7em}\rule{\linewidth}{0.4pt}\vspace{-0.4em}]
\titlespacing*{\section}{0pt}{1em}{0.4em}
\newenvironment{cvrole}[5]{%
  \par\noindent\textbf{\large #1} | \textit{\small #2, #3}\hfill \textit{\small #4 - #5}\par
  \begin{itemize}[leftmargin=2em, itemsep=-0.3em, topsep=-0.5em]%
}{%
  \end{itemize}\par\noindent\mbox{}\par\vspace{0.1em}%
}
\newcommand{\cvName}{Nishal Stanislaus Silva}
\newcommand{\cvEmail}{nishal.silva@hotmail.com}
\newcommand{\cvPhone}{+1 613 200 2030}
\newcommand{\cvCity}{Toronto, ON}
\newcommand{\cvSite}{https://nishal.xyz}
\newcommand{\cvLinkedIn}{https://www.linkedin.com/in/nishal-silva}
\newcommand{\cvGitHub}{https://github.com/ns2max}
\newcommand{\cvStatus}{Canadian Permanent Resident, Eligible to work in Canada without sponsorship}
\newcommand{\cvTagline}{Computer Vision \& ML Engineer | Industrial Inspection, IoT \& Edge Deployment}

\begin{document}
\pagestyle{empty}
\begin{center}
    {\LARGE \scshape \textbf{\cvName}}{\large \textit{, PhD}}\\[1pt]
    {\small \cvTagline}\\[1pt]
    {\footnotesize\textit{\cvStatus}}\\[1pt]
    {\small
    \href{mailto:\cvEmail}{\cvEmail} \,\,|\,\, \cvPhone \,\,|\,\, \cvCity
    \,\,|\,\, \href{\cvSite}{nishal.xyz} \,\,|\,\, \href{\cvLinkedIn}{linkedin.com/in/nishal-silva}
    \,\,|\,\, \href{\cvGitHub}{github.com/ns2max}}
\end{center}

\section*{Professional Summary}

Computer vision and ML engineer with a PhD (University of Trento, 2025) and 10+ years across industry and academia, focused on perception systems that ship on resource-constrained hardware. At MAS Holdings I built and deployed production computer-vision inspection for apparel manufacturing at scale: camera, sensor and IoT streams feeding OpenCV pipelines and embedded C++/Python inference into live workflows, cutting inspection time 99.5\% and lifting operational efficiency by up to 300\%. My doctoral work carried the same discipline to the model side: a real-time detector running inference in 14 ms on a Raspberry Pi 4 at 31.4\% CPU, F1 0.76, 74$\times$ faster than the DTW baseline it replaced (JAES 2026), plus a training-free method reaching 95\% detection with no training data. My work runs from data labeling and evaluation design through deployment, monitoring and CI/CD.

\section*{Core Competencies}

\textbf{Computer Vision \& Perception:} production OpenCV inspection pipelines, template and feature matching, image registration and geometric warping, corner and edge detection, OCR and layout segmentation, explainable defect overlays, industrial and microscopic camera arrays, multi-camera and multi-sensor capture, sensor fusion (IMU, audio, XR), CNN and sequence architectures (RNN, LSTM, GRU), open-set and similarity-based matching\\
\textbf{ML Engineering:} data-labeling and label-ontology design, training and evaluation under domain shift, benchmark and dataset design, ablation studies, evaluation-metric design tied to operational outcomes, deployment monitoring and failure diagnosis, TensorFlow/Keras, PyTorch, scikit-learn\\
\textbf{Systems \& Deployment:} real-time inference under sub-30 ms budgets, latency, memory, CPU and power profiling on ARM and Raspberry Pi, embedded Linux, C++17 real-time engines, PLC and conveyor integration, Docker, CI/CD, code review, pytest, reproducible pipelines, deployment documentation

\section*{Technical Stack}

\textbf{Languages:} Python, C++17, C, MATLAB, JavaScript/TypeScript, SQL, Shell\\
\textbf{ML \& AI:} TensorFlow, Keras, PyTorch, scikit-learn, NumPy, Pandas, SciPy; RNN/CNN, embedding and similarity matching, VAE and diffusion prototypes; familiarity with model compression, quantization, ONNX and TFLite; ARM latency profiling\\
\textbf{Computer Vision:} OpenCV, template and feature matching, OCR, image registration and warping, corner detection, explainable overlays, industrial and microscopic cameras, projector-camera AR\\
\textbf{Edge \& Systems:} Raspberry Pi (production), ARM CPU and latency profiling, embedded Linux, IMU/I\textsuperscript{2}C sensors, PLC interfacing, FFTW, libsndfile\\
\textbf{Data \& Dev:} AWS (EC2, S3, ECS, MWAA, RDS), ETL pipelines, Docker, Airflow, Jenkins, Git, Linux/UNIX, pytest, CI/CD, REST APIs

\section*{Education}
\textbf{Ph.D - Information Engineering and Computer Science} \textit{\small{ - University of Trento, Italy}} \hfill \textit{Jan 2025}\\[2pt]
\noindent\textit{\small Thesis: Embedded Real-time Musical Pattern Detection for Smart Musical Instruments}\\[3pt]
\noindent\textbf{M.Sc - Telecommunication and Electronic Engineering} \textit{\small{ - Sheffield Hallam University, UK}} \hfill \textit{Aug 2018}\\[3pt]
\noindent\textbf{B.Eng (Hons) - Electronic Engineering} \textit{\small{ - Sheffield Hallam University, UK}} \hfill \textit{Mar 2014}

\section*{Professional Experience}

\begin{cvrole}{Research Engineer}{MAS Holdings (Pvt) Ltd}{Colombo, Sri Lanka}{Jan 2015}{Jul 2020}
    \item Built and deployed end-to-end computer-vision and sensor systems for apparel manufacturing at scale, integrating camera, sensor and IoT streams into production workflows; operational efficiency improved by up to 300\%.
    \item Shipped automated multilingual care-label QC: OpenCV template and feature matching with an explainable defect overlay, iterated from scanner to industrial camera to a conveyor with PLC reject; inspection time cut 99.5\% and line operators could adjudicate borderline rejects without ML support.
    \item Built loom-side fabric defect detection with a microscopic camera array (overlapping fields of view across 60-inch rolls, 5 to 6 second alert budget); pilot cut operator headcount 50\%.
    \item Delivered the Promptly on-demand garment-printing system: designed the camera-to-warp geometry-alignment algorithm, personally built the RIP integration into the automation loop, and consulted on the flipping mechanism; now a Twinery/MAS commercial product across four countries.
    \item Built real-time ingestion and ETL for high-frequency IoT and operational time-series across sites in South Asia, North America and Africa; ran C++/Python inference on embedded hardware; introduced CI/CD and code review; mentored engineers and delivered group-wide CV training.
\end{cvrole}

\needspace{5\baselineskip}
\begin{cvrole}{Postdoctoral Researcher}{University of Trento}{Trento, Italy}{Jan 2025}{Dec 2025}
    \item Owned the full ML pipeline for streaming audio and IMU/gesture data: acquisition, feature engineering, data labeling and label-ontology design, training and evaluation under domain shift, edge deployment on Raspberry Pi 4 with monitoring; EU Horizon EIC Pathfinder project MUSMET.
    \item Delivered real-time inference at 14 ms on a Raspberry Pi 4, F1 0.76, 31.4\% CPU, 74$\times$ faster than the 1038 ms DTW baseline (JAES 2026); designed the frozen-protocol benchmark suites and ablations that quantified the accuracy, latency and CPU trade-offs.
\end{cvrole}

\needspace{5\baselineskip}
\begin{cvrole}{Visiting Researcher}{McGill University (IDMIL / CIRMMT)}{Montreal, Canada}{May 2024}{Aug 2024}
    \item Fused inertial (IMU) and audio streams through a hierarchical finite state machine, accumulating repeated observations of the same event to reach F1 0.78 across a 500-event corpus (IEEE IS\textsuperscript{2} 2025); profiled compute and power for real-time feasibility on self-powered embedded hardware.
\end{cvrole}

\section*{Selected Projects \footnotesize {\normalfont{(full portfolio: \href{https://nishal.xyz/research}{\textit{nishal.xyz/research}})}}}
\begin{itemize}[leftmargin=1.2em, itemsep=-0.25em]
    \item \textbf{Tech-pack digitization (MAS):} pre-LLM OCR and layout segmentation of physical garment blueprints into a structured, searchable archive.
    \item \textbf{Lace archive and ERP search (Noyon Dentelle):} flatbed and microscopic capture preserving relief, designer metadata tagging, context-aware search inside the ERP; still in production use.
    \item \textbf{Real-time audio pattern detection:} MFCC front-end into a stacked RNN on 30 ms windows; Raspberry Pi 4 at 14 ms, 31.4\% CPU, F1 0.76 against a DTW baseline at 1038 ms.
    \item \textbf{SSIM-based training-free detection:} four-attribute representation with bicubic resizing for variable length; 95\% detection across human, synthetic and JKUPDD sets with no training data.
\end{itemize}

\enlargethispage{4\baselineskip}
\vspace{-0.5em}
\section*{Highlights \& Recognition}
\begin{itemize}[leftmargin=1.2em, itemsep=-0.25em]
    \item \textbf{GMOA Recognition (Sri Lanka):} computer-vision device digitizing patient vitals from retrofitted webcams, deployed across hospital wards for COVID-19 remote monitoring.
    \item \textbf{Commercial impact:} Promptly on-demand printing, to which I contributed the alignment algorithm and RIP integration, is now deployed across facilities in four countries.
    \item \textbf{Research and open tooling:} 10+ peer-reviewed papers (JAES, IEEE, ACM); Nebula, a C++17 feature-extraction library with per-feature ARM latency benchmarks; four open-access Zenodo datasets; side projects built with AI-assisted workflows including Claude Code.
\end{itemize}

\end{document}
